   Dark energy and gravity are equally important in determining the large-scale geometry and evolution of the universe. This evolution can be computationally modelled using the Friedmann ODEs obtained from the general relativistic Friedmann-Robertson-Lemaitre-Walker metric tensor.  Under the assumption of a spatially flat universe with its only contents being matter and a time-varying dark energy, the parameters in the Friedmann equations reduce to the Hubble constant, the cosmological matter density, the two-parameter dark energy equation of state and a matter-dark energy interaction parameter. 

   I demonstrate the use of Bayesian inference to estimate values for these cosmological parameters (with the excpetion of the Hubble constant, as this cannot be extracted from the data). 

   The parameters are related to observations by redshift-distance relations. I focus on the luminosity distance, which can be measured using Type Ia supernovae. The data used is a selection of three recent SNe compilations and five artificially generated sets: three based on these compilations to see how the parameter uncertainty scales with error and two projected data sets based upon the upcoming Supernova Acceleration Probe. I use the Metropolis-Hastings algorithm, an example of Markov chain Monte Carlo methods, to generate the posterior probability distributions for models with one and two free parameters (although the code I wrote can be used for up to $n$ parameters). Credible regions can be generated for the parameter space from the posterior distributions. 

   I conclude that the one-parameter distributions show that the currently accepted ”concordance cosmology” is a plausible approximation to the parameters, although differences in the light curve fitters which were used to calibrate the data prevent a more decisive estimation. The two-parameter models exhibit a wide range of degeneracies which need to be resolved by combining the supernova-derived credible regions with those from the cosmic microwave background and baryon acoustic oscillations.

